\section{Результаты}
%------------------------------------------------
\begin{frame}
	\frametitle{Основные результаты работы}
	\begin{itemize}
		\item Изменение коэффициента проскальзывания в резонансной структуре повышает эффективность стохастического охлаждения, но эффекты ВПР делают предпочтительной регулярную структуру для экспериментов с тяжелыми ионами;
		\vspace{1em}
		\item Для коллайдерных экспериментов с поляризованными протонами резонансная структура повышает критическую энергию, искажая дисперсионную функцию;
		\vspace{1em}
		\item Численные исследования выявили нестабильность продольного фазового движения при прохождении критической энергии. Процедура скачка критической энергии помогает решить эту проблему. Экспериментальные данные с синхротрона У-70 согласуются с численными оценками, учитывающими высшие порядки разложения коэффициента уплотнения орбиты и импедансы для различных интенсивностей сгустка;
	\end{itemize}
\end{frame}
%------------------------------------------------
\begin{frame}
	\frametitle{Основные результаты работы}
	\begin{itemize}
		\item Скачок критической энергии ограничен его величиной и темпом изменения. Различия применения гармонического и барьерного ВЧ влияют на скачок. Оценки показывают ограничения на параметры сгустка из-за продольной микроволновой неустойчивости;
		\vspace{1em}
		\item Изучена спиновая динамика для ЭДМ. Предложена концепция квази-замороженного спина с обводными каналами. Фильтры Вина на прямых участках компенсируют поворот спина;
		\vspace{1em}
		\item Модернизированный синхротрон Nuclotron сохраняет функцию бустера для коллайдера NICA. В 8/16-периодичных структурах возможны эксперименты по ЭДМ и поиск аксиона.
	\end{itemize}
\end{frame}
%------------------------------------------------
\begin{frame}
	\frametitle{Апробация работы}
	Основные результаты работы были представлены докладывались~на:
	\small
	\begin{itemize}
		\item 63, 65, 66-ая Всероссийская научная конференция МФТИ в 2020, 2023, 2024 гг. г. Долгопрудный,
		Россия;
		\item XXVII и XXVIII Всероссийская конференции по ускорителям заряженных частиц RuPAC'21, RuPAC'23. Алушта; Новосибирск, Россия.
		\item VII, VIII, IX и X Международная конференция Лазерные и Плазменные технологии ЛаПлаз'21, ЛаПлаз'22, ЛаПлаз'23, ЛаПлаз'24, ЛаПлас'25. Москва, Россия;
		\item XIII, XIV, XVI международная конференция по ускорителям заряженных частиц IPAC'22 IPAC'23, IPAC'25. Бангкок, Тайланд; Венеция, Италия; Тайпей, Тайвань;
		\item XIX Международная конференции по спиновой физике высоких энергий DSPIN'23. Дубна, Россия;
		\item XI-я Международная конференция по ядерной физике в накопительных кольцах STORI’24. Хуэйчжоу, провинция Гуандун, Китай;
	\end{itemize}
\end{frame}
%------------------------------------------------
\begin{frame}
	\frametitle{Публикации}
	Основные результаты по теме диссертации изложены в 13 печатных
	изданиях: 10 печатных работ изданы в журналах, рекомендованных
	ВАК, 2 статьи — в журналах, индексируемых международными
	базами цитирования Scopus и Web of Science.
	\small
	\begin{itemize}
		\item	(to be published) Kolokolchikov, S.D., et al. Features of dual-purpose structure for heavy ion and light particles, Nucl.Sci. and Tech.
		\item (to be published) Kolokolchikov, S.D., et al. Quasi-frozen spin for both deuteron and proton beam at periodic EDM storage ring lattice, Nucl.Sci. and Tech.
		\item	Kolokolchikov, S.D., Senichev, Y.V. Magneto-Optical Structure of the NICA Collider with High Transition Energy. Phys. Atom. Nuclei 84, 1734–1742 (2021). https://doi.org/10.1134/S1063778821100185
		\item	Kolokolchikov, S., Senichev, Y. Peculiarities of Crossing and Raising the Synchrotron Transition Energy. Phys. Atom. Nuclei 86, 2260–2264 (2023). https://doi.org/10.1134/S1063778823110236
		\item	Kolokolchikov, S., Aksentyev, A., Melnikov, A. et al. Designing Bypass Channels in NICA Accelerator Complex for Polarized Beam Experiments for EDM Search. Phys. Atom. Nuclei 86, 2423–2428 (2023). https://doi.org/10.1134/S1063778823110248
		
	\end{itemize}
\end{frame}
%------------------------------------------------
\begin{frame}
	\frametitle{Публикации}
	\small
	\begin{itemize}
		\item	Kolokolchikov, S., Aksentiev, A., Melnikov, A. et al. Spin Coherence and Betatron Chromaticity of Deuteron Beam in “Quasi-Frozen” Spin Regime. Phys. Atom. Nuclei 86, 2684–2688 (2023). https://doi.org/10.1134/S106377882311025X
		\item	Acceleration and crossing of transition energy investigation using an RF structure of the barrier bucket type in the NICA accelerator complex / S. Kolokolchikov [et al.] // J.Phys.Conf.Ser. — 2023 — Vol. 2420, 012001 — DOI: 10.1088/1742-6596/2420/1/012001
		\item	Spin coherence and betatron chromaticity of deuteron beam in NICA storage ring / S. Kolokolchikov [et al.] // J.Phys.Conf.Ser. — 2024 — Vol. 2687, 022027 — DOI: 10.1088/1742-6596/2687/2/022027
		\item	ByPass optics design in NICA storage ring for experiment with polarized beams for EDM search / S. Kolokolchikov [et al.] // J.Phys.Conf.Ser. — 2024 — Vol. 2687, 022026 — DOI: 10.1088/1742-6596/2687/2/022026
	\end{itemize}
\end{frame}
%------------------------------------------------
\begin{frame}
	\frametitle{Публикации}
	\small
	\begin{itemize}
		\item	Kolokolchikov, S., Senichev, Y., Aksentiev, A. et al. Transition Energy Crossing of Polarized Proton Beam at NICA. Phys. Atom. Nuclei 87, 212–215 (2024). https://doi.org/10.1134/S1063778824700054
		\item	Kolokolchikov, S., Senichev, Y., Aksentev, A. et al. Longitudinal Dynamic in NICA Barrier Bucket RF System at Transition Energy Including Impedances in BLonD. Phys. Part. Nuclei Lett. 21, 419–424 (2024). https://doi.org/10.1134/S1547477124700389
		\item	Kolokolchikov, S.D., Senichev, Y.V. \& Kalinin, V.A. Transition Energy Crossing in Harmonic RF at Proton Synchrotron U-70. Phys. Atom. Nuclei 87, 1355–1362 (2024). https://doi.org/10.1134/S106377882410020X
		\item	Kolokolchikov, S.D., Aksentev, A.E., Mel’nikov, A.A. et al. Transition Energy Crossing in NICA Collider of Polarized Proton Beam in Harmonic and Barrier RF. Phys. Atom. Nuclei 87, 1449–1454 (2024). https://doi.org/10.1134/S1063778824100211
	\end{itemize}
\end{frame}
%------------------------------------------------


